% 本文件由 LaTeX 排版 Agent 根据有效 translation.json 自动生成。
% author=Venantius; work_id=0709; title=复活节之歌

\chapter{复活节之歌}

% structure: p0001 source=h1 target=omitted reason=page-title-already-rendered-by-parent

\Emph{\Person{文提乌斯·奥诺里乌斯·克莱门提亚努斯·福图纳图斯}的诗歌，\WorkTitle{论复活}}\par

季节因繁花与明媚天气而羞红，极地之门以更大的光明敞开。那喷火的太阳在天上的路径升得更高，他出发行走他的行程，进入海洋的水域。身披光线穿越流动的元素，在这短暂的夜晚，他把白昼伸展成一个圆圈。灿烂的天穹展露其清澈的面容，明亮的星辰显出它们的喜乐。大地以其多样的增益倾泻丰收的礼物，当年岁重新带来春日的丰盛。柔软紫罗兰的花床染紫了平原；草地因植物而翠绿，植物因叶片而闪耀。花朵的明亮光辉逐渐显现；百花以它们的绽放微笑。种子播下后，谷物在田野间远扬，许诺能战胜农夫的饥饿。藤蔓枝梢离弃了它的茎干，哀悼自己的欢乐；葡萄只从它惯常产出酒的水源里供水。含苞待放的嫩芽，从母亲背脊上带着柔软的茸毛升起，预备好它的胸膛以待生产。在冬季被撕裂了树叶后，翠绿的树林重新更新它的叶阴。柳树、冷杉、榛树、柳条、榆树、枫树、胡桃树，混杂在一起，每棵树都以其叶片欢欣鼓舞。于是蜜蜂，即将建造它的蜂巢，离开蜂箱，嗡鸣于花朵之上，用它的腿运走蜂蜜。那鸟儿，曾闭住歌声而沉默，因冬日的寒冷而迟钝，现在重归它的曲调。因此，\Person{菲洛墨拉}以自己的乐器调谐音调，空气也因回响的旋律而更加甜美。看哪，复苏世界的恩惠作证，所有礼物都与其主一同返回。\newline{}因为为着基督战胜后上升、降到阴郁的\Place{塔尔塔罗斯}的荣耀，树林到处以叶片表达赞同，植物以花朵表达赞同。光明、天空、田野和大海恰当地赞美那升到星辰之上、粉碎地狱律法的天主。看哪，那被钉十字架者作为天主治理万有，一切受造物向它们的造物主献上祈祷。万福，喜庆之日，遍受世人尊敬，在这一天，天主征服了地狱，赢得星辰！年岁和月亮的更替，白昼的丰盛之光，时辰的光辉，万物都以声音喝彩。因此，为着你的荣耀，树林以其叶片喝彩；因此，葡萄藤以其静默的嫩芽表达感恩。因此，灌木丛如今回响着鸟儿的低语；在这些之中，麻雀以充沛的情感歌唱。\newline{}哦，基督，世界的救主，仁慈的造物主与救赎者，惟一出自圣父天主性的后裔，以无可言喻的方式从你父的心中流出，你是自存的圣言，从你父亲的口中发出大能，与他同等，与他同心，与他同伙，与圣父同永，万物最初从他那里获得起源！你悬起穹苍，你堆聚土壤，你倾泻海洋，靠着你的治理，一切固定于其位置的事物繁荣兴旺。你看见人类沉沦于苦难的深处，为要拯救人，你自身也成了人：你不仅愿意有身体而生，而且成了血肉，这血肉忍受了出生和死亡。你经历葬礼，你自己是生命的作者和世界的塑造者，你踏上了死亡之路，赐予救恩的扶助。地狱律法的阴郁锁链屈服了，\OriginalTerm{混沌}{chaos}惧怕被光明之临在压迫。黑暗消逝，被基督的光辉驱散；永恒之夜的浓厚黑幕落下。\newline{}但请归还那应许的抵押，我祈祷，哦，良善的大能！第三天已到来；起来，我埋葬的主；你的肢体不应躺在低下的坟墓里，也无价值的石头应压住那世界的赎价。一块石头用狭窄的岩石封住并覆盖他，这不相称，在他手中万物都包容。请拿走细麻布，我祈祷；把汗巾留在坟墓里：你为我们足够，没有你便一无所有。释放地狱牢狱中被锁的阴影，召回那沉到最低深渊的一切。归还你的面容，使世界得见光明；归还那因你的死亡而逃离我们的白日。\newline{}但返回时，哦，神圣的征服者！你完全充满了天堂！\Place{塔尔塔罗斯}被压下，没有保有它的权利。下界的统治者，贪婪地张开他空洞的下颚，一直是掠夺者，却成了你的猎物。你从死亡的牢狱中救出无数人民，他们自由地跟随，随他们的领袖所到之处。猛兽惊慌中吐出它吞下的群众，羔羊从豺狼之口夺回羊群。因此，从下界重新寻找坟墓，你重取血肉，如同战士将丰硕的战利品带回天上。混沌曾因惩罚而拘留的，如今他已归还；死亡可能追寻的，新生命却持守着。\newline{}哦，神圣的君王，看哪，你胜利的极大一部分闪耀显现，当神圣的圣洗祝福纯洁的灵魂！一队披白衣的军团从明亮的水波出来，在全新的水流中洗净他们旧日的过犯。白衣也表明明亮的灵魂，牧者从雪白的羊群获得喜悦。\Person{费利克斯}司铎加入分享这奖赏，他愿意将双倍的塔冷通交与他的主。吸引那些在外邦人的错误中徘徊的人趋向更善，以免野兽将他们掳走，他守护天主的羊栈。那些先前被有罪的厄娃感染的人，如今在教会怀抱中，以丰盛之乳喂养，重新得救。以温和的谈话培育质朴的乡村心田，从荆棘中，因\Person{费利克斯}的恩惠，产生了庄稼。撒克逊人，一个凶猛民族，仿佛按野兽方式生活，当你，哦，圣者！施以医治时，野兽便相似绵羊。他们将要与你共度一个时代，以百倍的收获，你以丰收之产充满谷仓。愿这百姓，无玷污，在你手中坚固，愿你将纯洁的抵押带到星辰之上。愿一顶冠冕从高天赐予你，得自你自身，愿另一顶冠冕因你的百姓而繁茂。\par

\SourceRecord{Translated by William Fletcher.；Ante-Nicene Fathers；Vol. 7.；Edited by Alexander Roberts, James Donaldson, and A. Cleveland Coxe.；Buffalo, NY: Christian Literature Publishing Co.,；1886.；<http://www.newadvent.org/fathers/0709.htm>.}
